% !TeX root = ../main.tex
\section{Intensity Sweep and the Throughput--Latency Knee}\label{sec:intensity}
\rubric{Throughput--latency knee}{35}{Clear knee identification, diminishing-returns discussion and a
Little's Law connection.}

\begin{hypothesis}
  \todo{Thread count at which throughput should stop scaling and latency should rise, and why
  (cores, memory channels, shared bandwidth).}
\end{hypothesis}

\begin{knobs}
  Fixed       & region \tbd~per thread; random reads; granularity \tbd~bytes \\
  Varied      & threads $T$ = 1, 2, 4, \ldots{} (pinned: \todo{CPU list}) \\
  Script      & \code{Cache_and_Memory_Performance_Profiling/scripts/run.sh intensity} \\
\end{knobs}

\begin{figure}[tbp]
  \centering
  \IfDataExists{intensity}{%
    \begin{tikzpicture}
      \begin{axis}[acs, width=0.75\linewidth, xlabel={Throughput (GiB/s)}, ylabel={Latency (ns/op)},
                   legend pos=north west, nodes near coords, point meta=explicit symbolic,
                   every node near coord/.append style={font=\scriptsize, anchor=south east}]
        \addplot+ table[x=gib_per_s_median, y=lat_mean_ns_median, meta=threads] {data/intensity.csv};
        \addlegendentry{mean}
        \addplot+ table[x=gib_per_s_median, y=lat_p99_ns_median, meta=threads] {data/intensity.csv};
        \addlegendentry{p99}
      \end{axis}
    \end{tikzpicture}%
  }
  \caption{Throughput vs.\ latency; labels give the thread count. \todo{mark the knee}}\label{fig:intensity}
\end{figure}

\subsection{Little's Law check}
With $T$ threads always busy, the number of requests in the system is $L = T$, so throughput $\lambda$
and mean latency $W$ should satisfy $L = \lambda W$ \cite{little1961proof}.

\begin{table}[tbp]
  \centering\small
  \caption{Little's Law: measured $\lambda W$ against the number of busy threads.}\label{tab:littles-law}
  \IfDataExists{intensity}{%
    \pgfplotstabletypeset[
      col sep=comma,
      columns={threads,ops_per_s_median,lat_mean_ns_median,lambdaW},
      create on use/lambdaW/.style={create col/expr={\thisrow{ops_per_s_median}*\thisrow{lat_mean_ns_median}/1e9}},
      columns/threads/.style={column name=$T$},
      columns/ops_per_s_median/.style={column name={$\lambda$ (ops/s)}, sci, precision=2},
      columns/lat_mean_ns_median/.style={column name={$W$ (ns)}, fixed, precision=2},
      columns/lambdaW/.style={column name={$L=\lambda W$}, fixed, precision=2},
      every head row/.style={before row=\toprule, after row=\midrule},
      every last row/.style={after row=\bottomrule},
    ]{data/intensity.csv}%
  }
\end{table}

\subsection{Analysis}
\todo{Where the knee is and why; diminishing returns beyond it; whether $\lambda W \approx T$ holds and
what a deviation means (sampling, think time, oversubscription).}
